\chapter{Formulaire}

\noindent\textit{Toutes les formules du programme de 3ème, regroupées par thème pour une
révision rapide. Chaque formule est démontrée et illustrée dans le chapitre correspondant :
ce formulaire n'est qu'un aide-mémoire, pas un cours.}
\vspace{8pt}

\begin{multicols}{2}

\subsection*{\color{onbo}La mole et quantité de matière}
\begin{itemize}
\item Quantité de matière : $n=\dfrac mM$ ($m$ en g, $M$ en g/mol, $n$ en mol)
\item Nombre d'entités : $N=n\times N_A$, avec $N_A\approx6{,}02\times10^{23}$ mol$^{-1}$
\item Masse molaire moléculaire : somme des masses molaires atomiques des atomes de la
molécule
\end{itemize}

\subsection*{\color{onbo}Réactions chimiques}
\begin{itemize}
\item Loi de Lavoisier : masse des réactifs $=$ masse des produits (conservation de la
masse et des atomes)
\item Équilibrer une équation : ajuster les coefficients pour retrouver le même nombre de
chaque type d'atome de part et d'autre
\end{itemize}

\subsection*{\color{onbo}Électrolyse et solutions aqueuses}
\begin{itemize}
\item Électrolyse de l'eau : $2\,\mathrm{H_2O}\longrightarrow2\,\mathrm{H_2}+\mathrm{O_2}$
(rapport des volumes $\mathrm{H_2}:\mathrm{O_2}=2:1$)
\item Synthèse de l'eau : $2\,\mathrm{H_2}+\mathrm{O_2}\longrightarrow2\,\mathrm{H_2O}$
\item Concentration molaire : $c=\dfrac nV$ ($n$ en mol, $V$ en L, $c$ en mol/L)
\item Échelle de pH : de $0$ (acide) à $14$ (basique), $7$ = neutre
\end{itemize}

\subsection*{\color{onbo}Machines simples}
\begin{itemize}
\item Poids : $P=m\times g$, avec $g\approx10$ N/kg
\item Plan incliné : $F=\dfrac{P\times h}{L}$
\item Poulie fixe : $F=P$ \quad Poulie mobile : $F=\dfrac P2$ \quad Palan à $n$ brins :
$F=\dfrac Pn$
\item Treuil : $F=\dfrac{P\times r}{R}$ ($r$ = rayon du tambour, $R$ = longueur de la
manivelle)
\end{itemize}

\subsection*{\color{onbo}Loi d'Ohm et circuits électriques}
\begin{itemize}
\item Loi d'Ohm : $U=R\times I$
\item Circuit série : $I$ commun ; $U=U_1+U_2$ ; $R_{\text{tot}}=R_1+R_2$
\item Circuit en dérivation : $U$ commun ; $I=I_1+I_2$ ;
$\dfrac1{R_{\text{tot}}}=\dfrac1{R_1}+\dfrac1{R_2}$
\item Puissance électrique : $P=U\times I$
\end{itemize}

\subsection*{\color{onbo}Production et transport de l'énergie électrique}
\begin{itemize}
\item Pertes en ligne (effet Joule) : $P_{\text{pertes}}=R\times I^2$
\item Transformateur : $\dfrac{U_2}{U_1}=\dfrac{n_2}{n_1}$ ; conservation de puissance
(idéale) : $U_1I_1=U_2I_2$
\end{itemize}

\subsection*{\color{onbo}Tension alternative et énergie domestique}
\begin{itemize}
\item Fréquence et période : $f=\dfrac1T$ (secteur : $f=50$ Hz, $T=0{,}02$ s)
\item Énergie électrique : $E=P\times t$ ($1$ kWh $=1\,000$ Wh)
\item Coût d'une consommation : coût $=$ énergie (kWh) $\times$ prix unitaire du kWh
\end{itemize}

\subsection*{\color{onbo}Pétrole et combustion}
\begin{itemize}
\item Combustion complète : hydrocarbure $+$ dioxygène (suffisant) $\to\mathrm{CO_2}+
\mathrm{H_2O}$
\item Combustion incomplète : dioxygène insuffisant $\to$ production de CO (toxique)
\item Ordre des fractions du pétrole (léger $\to$ lourd) : gaz, essence, kérosène, gasoil,
fioul lourd, bitume
\end{itemize}

\subsection*{\color{onbo}Matières plastiques}
\begin{itemize}
\item Polymérisation : $n\,\text{monomères}\longrightarrow\text{polymère}$ (ex. :
$n\,\mathrm{CH_2\!=\!CH_2}\to\mathrm{(-CH_2-CH_2-)}_n$)
\item Thermoplastique : ramollit à la chaleur, refondable. Thermodurcissable : durcit
définitivement
\end{itemize}

\subsection*{\color{onbo}Dessin technique}
\begin{itemize}
\item Échelle : $\text{échelle}=\dfrac{\text{dimension dessin}}{\text{dimension réelle}}$
\item $1{:}n$ = réduction ; $n{:}1$ = agrandissement ; les cotes indiquées sont toujours
les dimensions \textbf{réelles}
\item Perspective cavalière : face en vraie grandeur ; fuyantes à $45°$, coefficient
souvent $0{,}5$
\end{itemize}

\subsection*{\color{onbo}Transmission du mouvement de rotation}
\begin{itemize}
\item Engrenage : $N_1\times Z_1=N_2\times Z_2$
\item Poulies-courroie : $N_1\times D_1=N_2\times D_2$
\item Rapport de transmission : $r=\dfrac{N_2}{N_1}=\dfrac{Z_1}{Z_2}$ (ou
$\dfrac{D_1}{D_2}$) ; $r<1$ réducteur, $r>1$ multiplicateur
\end{itemize}

\subsection*{\color{onbo}Moteurs}
\begin{itemize}
\item Cycle à $4$ temps : admission $\to$ compression $\to$ explosion (temps moteur)
$\to$ échappement
\item Rendement : $\eta=\dfrac{E_{\text{utile}}}{E_{\text{fournie}}}\times100$
\end{itemize}

\subsection*{\color{onbo}Installations électriques domestiques}
\begin{itemize}
\item Puissance maximale admissible : $P_{\max}=U\times I_{\text{calibre}}$
\item Disjoncteur : protège le \textbf{circuit} (surintensité), réarmable
\item Différentiel $+$ mise à la terre : protègent les \textbf{personnes} (fuite de
courant)
\end{itemize}

\end{multicols}
